% JAC1964-20261121, /root/arithmetic_sign. Research reduction; no positivity closure.
\subsection*{AR-29. Full-Jacobian diagonal terms and their remaining constraints}

The Jacobian cone in AR-28 is larger than the fixed $E^3$ cone in
AR-26.  The following exact decompositions identify, rather than
silently discard, the additional terms.  They do not establish a sign
for the full complementary group.

\paragraph{AR-29a: the quotient and its orthogonal complement.}
Let $C/K,e,f:C\to A$ be a parametrized test curve as in AR-27,
and write $J=\operatorname{Jac}(C)$.  Give $J$ its principal
polarization and $A=E^3$ the product polarization.  The map $f$
induces a surjective homomorphism $F:J\to A$.  Its polarized adjoint
is denoted $F^\dagger$.  The positive endomorphism
$S=FF^\dagger$ of $A$ is an isogeny.  On the rational Hodge realizations, $F$ is
surjective and its polarized adjoint is injective; the positive
Hermitian forms therefore make $FF^\dagger$ invertible.  A
homomorphism of abelian varieties inducing an isomorphism on
$H^1$ is an isogeny.  This is also the usual nondegeneracy of
the polarization induced on an abelian subvariety.  Hence
\[
 p=F^\dagger S^{-1}F,\qquad q=1-p
\]
are complementary Rosati-self-adjoint rational idempotents on $J$.
Let $B$ be the abelian subvariety up to isogeny defined by $q$.
There is an exact Chow-motive decomposition
\[
 h^1(C)=p h^1(C)\oplus q h^1(C)
       \simeq h^1(E^3)\oplus h^1(B).
\]
Here the correspondence construction on $h^1$ is fully faithful
for homomorphisms of abelian varieties.  Thus these are exact
endomorphism and Chow identities; no equality of projectors is
being inferred just from realization.  The word ``complement''
does not require $B$ to have no further factors isogenous to $E$.

Let $\widehat t_R$ denote the uniquely admissibly metrized,
axes-normalized correspondence attached to $R\in\operatorname{End}^0(J)$.
In particular $t_1=t_e$ is the normalized diagonal, and
$\widehat t_e=\widehat t_p+\widehat t_q$ as adelic line bundles.
For $R_1,R_2$ supported on $p$, Zhang's established tensor identity
therefore gives
\[
 \begin{split}
 \langle\widehat t_{R_1}\widehat t_{R_2}\widehat t_e\rangle
 &=\langle\widehat t_{R_1}\widehat t_{R_2}\widehat t_p\rangle
   +\langle\widehat t_{R_1}\widehat t_{R_2}\widehat t_q\rangle,\\
 \langle\widehat t_{R_1}\widehat t_{R_2}\widehat t_q\rangle
 &=h\bigl(\Delta_e,(R_1\otimes R_2\otimes q)\Delta_e\bigr),
 \end{split}                                                    \tag{AR.29.1}
\]
up to the harmless simultaneous transpose if the pullback convention
is used.  All heights are divided by $[K:\mathbf Q]$.
The second term is a genuine additional term; its vanishing is
not implied by primitivity of the target class in $\operatorname{CH}^2(N)$.

\paragraph{AR-29b: the precise higher-level part of the extra term.}
On two selected elliptic factors let $u=[i]$ act on $h^1(E)$, so
$u^2=-1$.  The exact self-transpose idempotents
\[
 P_{\rm same}=\frac12(1-u\otimes u),\qquad
 P_{\rm opp}=\frac12(1+u\otimes u)                       \tag{AR.29.2}
\]
split $h^1(E)^{\otimes2}$.  After adjoining $i$ to the
coefficients they select respectively equal and opposite CM
characters.  The opposite-character motive is
$\mathbf Q(-1)^{\oplus2}$; its two divisor classes may be taken
as the normalized diagonal and the normalized graph of $[i]$ on
$E^2$.  Their intersection matrix is $-2I_2$ (replace the second
class by its negative if the graph convention requires it).
Indeed the squares are $-2$, while their mixed intersection is
$\deg(1-i)-2=0$; these two independent divisor classes exhaust
the rank-two opposite-character realization.  The inverse
intersection matrix supplies a Chow splitting.  Since the
resulting divisor projector commutes with (AR.29.2), the remaining
homologically trivial idempotent vanishes by Kimura nilpotence.
Thus the displayed Tate identification is at the Chow level.

Consequently the extra motive in (AR.29.1) has a direct summand
\[
 P_{\rm opp}h^1(E)^{\otimes2}\otimes h^1(B)
      \simeq h^1(B)(-1)^{\oplus2}.                     \tag{AR.29.3}
\]
Self-pairings on this block are controlled by the ordinary
divisor height pairing.
For example, realizing $h^1(B)$ as $q h^1(C)$ and writing the
corresponding cycles as $D_0\times\beta_0+D_1\times\beta_1$
on $E^2\times C$, with $\beta_j\in q\operatorname{Pic}^0(C)_\mathbf Q$,
height adjunction and $D_iD_j=-2\delta_{ij}$ give
\[
 h(D_0\times\beta_0+D_1\times\beta_1,
       D_0\times\beta_0+D_1\times\beta_1)
       =4\bigl(\widehat h_{\theta_J}(\beta_0)
                    +\widehat h_{\theta_J}(\beta_1)\bigr)\ge0,   \tag{AR.29.4}
\]
where the standard divisor convention is
$h_C(\beta,\beta)=-2\widehat h_{\theta_J}(\beta)$.
For two finite elliptic quotient maps $f_j:C\to E$ of degrees
$d_j$, transport by $f_1^*\otimes f_2^*\otimes q$ multiplies
these self-heights by $d_1d_2>0$, since its transpose composed
with itself is $d_1d_2$ on the selected source motive.  This
observation concerns self-pairings; it does not assign a sign to
an arbitrary mixed pairing in (AR.29.1).

In contrast, put $K=P_{\rm same}h^1(E)^{\otimes2}$.
It has Hodge types $(2,0)$ and $(0,2)$.  If $b=\dim B>0$, the
Hodge numbers of $K\otimes h^1(B)$ are
\[
 h^{3,0}=h^{2,1}=h^{1,2}=h^{0,3}=b.                    \tag{AR.29.5}
\]
In particular the whole motive cannot be replaced by a Tate twist
of an abelian $h^1$ motive, since such a twist has only types
$(2,1)$ and $(1,2)$.  Formula (AR.29.5) neither proves that a given
projected diagonal class is nonzero nor rules out a further
level-one direct summand.  It does show exactly why the extra
term is not entirely covered by Neron--Tate positivity.
No source checked here gives a sign for this same-character part
on every auxiliary curve.

\paragraph{AR-29c: the new Prym contributes no one-Prym term.}
For the double cover $\rho:C'\to C$ of AR-27 let $\iota$ be its
deck involution.  It fixes the ramification point $e'$ and hence
\[
 (\iota,\iota,\iota)_*\Delta_{e'}=\Delta_{e'}
\]
term by term.  On $h^1(C')$ the exact old and Prym projectors are
\[
 p_0=\tfrac12\rho^*\rho_*=\tfrac12(1+\iota_*),
 \qquad q_0=\tfrac12(1-\iota_*).
\]
Therefore
\[
 (p_0\otimes p_0\otimes q_0)_*\Delta_{e'}=0,           \tag{AR.29.6}
\]
and similarly after permuting the factors: the left side is
simultaneously invariant and anti-invariant under the diagonal
deck involution.  Thus canonicalization creates no additional
mixed term with two old elliptic factors and exactly one new
Prym factor.  An original complement $B\subset J(C)$ is not
removed by this observation.

There is a useful exact normalization check.  For normalized
correspondences $t_1,t_2$ on $C$, put
$t'_j=\rho^*t_j\rho_*$.  Their canonical metrics pull back under
$\rho\times\rho$, by rigidification and functoriality.  Using
(AR.29.6), $r_*\Delta_{e'}=2\Delta_e$ and $r=\rho^3$, one obtains
\[
 \begin{split}
 \langle\widehat t'_1\widehat t'_2\widehat t_{e'}\rangle
 &=h\bigl(\Delta_{e'},(t'_1\otimes t'_2\otimes p_0)
                                    \Delta_{e'}\bigr)\\
 &=\tfrac12 h\bigl(r_*\Delta_{e'},
                         (t_1\otimes t_2\otimes1)r_*\Delta_{e'}\bigr)\\
 &=2\langle\widehat t_1\widehat t_2\widehat t_e\rangle.
 \end{split}                                                    \tag{AR.29.7}
\]
Three pulled-back cross correspondences instead give the factor
four by the degree-four map $\rho\times\rho:C'^2\to C^2$.
Neither calculation asserts that the vertical terms in a nef
polarization vanish.  Those terms still occur in the divisor
Hodge-index inequalities, so (AR.29.6)--(AR.29.7) alone do not orient the
remaining cubic.

\paragraph{Source boundary.}
The metric linearity and tensor-height identity in (AR.29.1) and
(AR.29.7) use Zhang 2010, Lemmas 2.2.1--2.2.4 and Theorem 2.3.5,
pp.~20--27, directly reread.  Its Conjecture 2.4.1, p.~28, is
still the general projected-positivity statement.  The divisor
normalization in (AR.29.4) is the Faltings--Hriljac formula, also
stated in Yu--Yuan--Zhou 2026, Theorem 2.5, p.~32.  The direct
sum, CM character calculations and deck-involution cancellation
above are deductions, not new source assumptions.  This note
explains the extra terms in the larger Jacobian cone; it does
not claim that its entire constraint system has already been
reduced to AR-26.

\paragraph{AR-29d: scalar verticals give a larger exact symmetry.}
Keep the actual canonical point $e$ and the metrics of AR-28.
Let $\overline\alpha=\overline\omega_a/(2g-2)$ and
$\overline c=\langle\overline\alpha,\overline\alpha\rangle$.
For any normalized cross bundle $\overline Q_R$ on $C^2$ induced
by an endomorphism $R$ of $J$, relative Deligne pairing gives
\[
 \langle\overline\alpha_1,\overline Q_R\rangle_{C^2/C}=0,
 \qquad
 \langle\overline\alpha_2,\overline Q_R\rangle_{C^2/C}=0.       \tag{AR.29.8}
\]
The projections in the two expressions integrate the factor
containing $\alpha$.  Indeed, AR-28 gives
$\overline\alpha=\overline{\mathcal O(e)}_a+\lambda\pi^*\overline c$
with $\lambda=(2g-1)/(2g)$.  The base term pairs to zero because
$Q_R$ has relative degree zero.  The remaining relative pairing
is the restriction of $Q_R$ to the section $e$, which is
metrically trivial by axis rigidification.  For the equality with
that restriction, the induction formula has no integral error:
$Q_R$ is an admissible degree-zero bundle on every curve fiber.
This argument proves identities of adelic line bundles, not only
zero degrees.

It follows that every triple intersection with exactly one cross
factor and two factors from
$\langle\alpha_1,\alpha_2,\pi^*c\rangle$ vanishes.  Put
\[
 A_j=g\alpha_j-\tfrac12\pi^*c=i_e^*\theta_J\text{ on factor }j.
\]
Consider all canonical nef pullbacks from $J^2$ of the form
\[
 a A_1+b A_2+Q_R+s\pi^*c,
 \qquad
 \begin{pmatrix}aI&R\\R^\dagger&bI\end{pmatrix}\succeq0,
 \quad s\ge0.                                             \tag{AR.29.9}
\]
The involution fixing $\alpha_1,\alpha_2,c$ and replacing every
$Q_R$ by $-Q_R$ is well defined by axis normalization and the
injectivity of $\operatorname{End}^0(J)\to\operatorname{Pic}^-(C^2)$.
It preserves this cone, by conjugation with $\operatorname{diag}(I,-I)$,
and preserves the geometric intersection form.  It fixes all
arithmetic terms with an even number of cross factors and reverses
all three-cross terms, by (AR.29.8).
Consequently the entire divisor Hodge-index/nef-triple constraint
system on this space, with its geometric and even arithmetic data,
is unchanged under that transformation of the intersection
functional.  Every AR-28 nef divisor belongs to (AR.29.9), with
its exact nonnegative base coefficient as in (AR.28b).  Thus this
statement does cover those extra canonical-diagonal nef tests.
It does not include every canonically metrized vertical pullback
of an arbitrary polarization on $J$.

This last symmetry does \emph{not} preserve independently specified
positive self-height values of all Picard summands: it reverses
those cubics too.  It is therefore an obstruction only for the
stated Hodge-index/nef constraints plus their even data.  It is
not by itself an obstruction after adjoining all Neron--Tate
information.

\paragraph{AR-29e: reflection preserves the two-$A$/one-$B$ data.}
There is a more selective statement for the narrow diagonal test
family.  Set $S=-p+q$, a Rosati-unitary rational involution, and
restrict cross operators to the block-diagonal algebra
\[
 \mathcal R=p\operatorname{End}^0(J)p
                 \oplus q\operatorname{End}^0(J)q.
\]
For $R=R_A+R_B$ in this algebra define
$\sigma_S(Q_R)=Q_{SR}=-Q_{R_A}+Q_{R_B}$ and fix the scalar
verticals and $c$.  The cone (AR.29.9) restricted to
$R\in\mathcal R$ is preserved, by conjugation with
$\operatorname{diag}(S,I)$.

The even arithmetic data are preserved as well.  The only point
beyond (AR.29.8) is
\[
 \langle\alpha_j,Q_{R_A},Q_{R_B}\rangle=0,
 \qquad j=1,2.                                          \tag{AR.29.10}
\]
To prove it, compose the two admissible normalized correspondences
along one curve factor.  Their composition is zero because
$R_A^\dagger R_B=R_A R_B^\dagger=0$.
Zhang's composition law, defined by relative Deligne pairing on
$C^3$, respects the admissible metrics.  Pulling this zero
composition to the diagonal in the two remaining variables gives
$\langle Q_{R_A},Q_{R_B}\rangle_{C^2/C}=0$ for either projection.
The Deligne-pairing projection formula gives (AR.29.10).
The same proof includes base factors, while geometric pairings
are preserved by Rosati orthogonality.

Thus the reflection changes a three-cross coefficient with $k$
$A$-supported factors by $(-1)^k$.  In particular it reverses the
$A^3$ coefficient containing the target $N$ height, while fixing
every $A^2B$ coefficient, including all of the positive
Tate--Picard self-heights in (AR.29.4).  This proves a symmetry
of the scalar-vertical, block-diagonal divisor constraint system
that retains this entire additional two-$A$/one-$B$ data.
It applies in particular to tests whose primitive divisor has an
$A$-supported cross part and whose nef polarization is taken from
the canonical diagonal family of AR-28, with any additional
$A$-supported cross terms permitted by the same cone.

The restriction to block-diagonal cross operators is substantive.
For an arbitrary off-diagonal $R$, left multiplication by $S$
still preserves the positive semidefinite cone, but need not fix
all arithmetic two-cross/vertical coefficients.  Those are not
being claimed invariant here.  Similarly, a one-$A$/two-$B$
Picard coefficient would change sign.  The next calculation kills
the canonical Lefschetz contraction in that sector, but not every
possible extra Picard summand of $h^1(A)\otimes h^2(B)$.

\paragraph{AR-29f: the canonical one-$A$/two-$B$ contraction is zero.}
Let $L_p,L_q$ be the rational symmetric lines on $J$ corresponding
to $p,q$ under the principal polarization, so
$L_p+L_q=\theta_J$.  Write
\[
 \beta_R=[i_e^*L_R-\deg(i_e^*L_R)e]\in J(K)_\mathbf Q.
\]
The canonical point condition and (AR.28c) give
$\beta_p+\beta_q=\beta_1=0$.
There is a symmetric rational line $L$ on $A$ such that
$L_p=F^*L$: its polarization operator is $(FF^\dagger)^{-1}$.
By AR-26, every symmetric divisor on $A$ kills the target
$Z=\operatorname{Ce}(f_*C)$ in Chow.  Its Albanese consequence is
\[
 F\beta_p=\operatorname{alb}_A(L\cdot f_*C)=0.
\]
Here the subtraction of the degree times $e$ contributes zero
because $f(e)=0$, and the even part of $f_*C$ contributes zero
to the Albanese, exactly as in AR-26.  Therefore
\[
 p\beta_p=F^\dagger(FF^\dagger)^{-1}F\beta_p=0,
 \qquad p\beta_q=0.                                    \tag{AR.29.11}
\]
This is the desired vanishing of the Lefschetz contraction.
To check its cycle normalization, let $t_q=-Q_q$ be the
normalized degree-zero correspondence associated with $q$.
Contraction of the first two factors of $\Delta_e$ with $t_q$
and pushforward to the third gives
\[
 \operatorname{pr}_{3*}(\operatorname{pr}_{12}^*t_q\cdot\Delta_e)
       =-2\beta_q.                                    \tag{AR.29.12}
\]
Indeed the small diagonal contributes
$\Delta^*t_q=-2i_e^*L_q$, the partial diagonal with its third
coordinate $e$ subtracts its degree times $e$, and the other
terms restrict $t_q$ to a rigidified axis and vanish.  Applying
$p$ to (AR.29.12) gives zero by (AR.29.11).
For a different algebraic class in $h^2(B)$, (AR.29.11) supplies
no such conclusion.  No positivity or vanishing for those
additional summands is assumed.

\paragraph{Status of AR-29.}
Parts (a)--(c) identify the extra motives and eliminate the new
Prym term exactly.  Parts (d)--(f) close the specified
canonical-diagonal Hodge-index route at the level of its
constraints, even while retaining the extra two-$A$/one-$B$
Picard self-heights and the zero Lefschetz contraction.
They neither prove the full arithmetic height conjecture nor
construct a negative self-height.  The transformed functional
comes from differing maps on the two curve factors; it does not
represent a new negative self-pairing of the same cycle.
Any use of further off-diagonal data or further Picard summands
must be checked separately.

\paragraph{AR-29g: a canonical point alone does not kill the Picard term.}
There is an explicit control example for this potential vanishing claim.
Let
\[
 C_0:\ y^2=(x^2-1)(x^2-2)(x^2-3),\qquad
 E:\ v^2=u^3-u,
\]
and take the smooth projective models.  The polynomial on the
right has six distinct roots, so $C_0$ has genus two.  The map
$f_0(x,y)=(x^2-2,y)$ is a degree-two morphism to $E$.
Let $\xi=(1,0)$, a Weierstrass point, and translate $f_0$ by
its two-torsion value at $\xi$ to obtain $F_0(\xi)=0$.
Let $q$ be the Rosati-orthogonal complement to the elliptic
quotient in $\operatorname{Jac}(C_0)$; its image $B$ has
dimension one.  Clear a denominator in $q$ when forming a
morphism of varieties; this does not affect any non-torsion claim.
The resulting morphism $q i_\xi:C_0\to B$ is nonconstant,
because $i_\xi(C_0)$ generates the Jacobian and $q\ne0$.
Choose an algebraic point $e$ such that
$q[e-\xi]$ is non-torsion.  Such a point exists: an elliptic
curve over a number field has algebraic points of arbitrarily
large canonical height (choose unbounded algebraic Weierstrass
coordinates and use $\widehat h=\tfrac12 h(x)+O(1)$), and the
nonconstant morphism $C_0\to B$ is surjective.  All these data
are defined after one finite extension of the number field.

Put $F=F_0-F_0(e)$, so $F(e)=0$.  Since $\xi$ is a
Weierstrass point, the hyperelliptic involution acts as $-1$ on
$\operatorname{Jac}(C_0)$ and satisfies $F_0\iota=-F_0$.
The divisor $F_0^*[0]-2\xi$ is invariant under this involution,
so its degree-zero Picard class is two-torsion and vanishes with
rational coefficients.  Moreover
$F^*[0]-F_0^*[0]$ lies in the image of $F_0^*\operatorname{Pic}^0(E)$,
which $q$ kills.  Hence
\[
 q[F^*[0]-2e]=-2q[e-\xi]\ne0.                         \tag{AR.29.13}
\]
Let $D_0=\Delta_E-E\times[0]-[0]\times E$.
Its restriction under $(F,F)$ is $-2F^*[0]$, since the normal
bundle of the elliptic diagonal is trivial.  Therefore its
normalized diagonal contraction has complementary Picard class
\[
 q\bigl[(F,F)^*D_0-\deg((F,F)^*D_0)e\bigr]
     =4q[e-\xi]\ne0.                                 \tag{AR.29.14}
\]
The class $D_0$ belongs to the opposite-character Tate block
of (AR.29.2).

Apply AR-27 to $(C_0,e)$, obtaining $\rho:C'\to C_0$ and an
actual canonical point $e'$ with $\rho^*e=2e'$.
Use $F\rho$ for both elliptic maps.  The class in (AR.29.14)
pulls back to its nonzero image under $\rho^*$; injectivity
follows from $\rho_*\rho^*=2$.
Thus there is a pointed curve with $K_{C'}\sim(2g(C')-2)e'$
whose two-elliptic-factor/opposite-character Picard contraction
is nonzero.  The corresponding Tate--Picard summand of the
projected threefold cycle has positive self-height by (AR.29.4).
This example tests the canonical-point assertion alone: its two
elliptic maps coincide, and it is not asserted to be a full
three-quotient AR-25 presentation.  It is not a negative-height
example.  For an AR-25 presentation the actual extra Picard
classes remain those in (AR.29.1), and their vanishing must be
proved from further information if it is needed.
